\NSPage{043}%
\begin{EESegment}{EE-TGT-P043-001}
and their inverses depend smoothly on \NSi{NS-F-P043-001}{\eta}. The inverse bounds are finite
because \NSi{NS-F-P043-002}{\lambda}, the radii, and all the bump functions were fixed before
\NSi{NS-F-P043-003}{N}.
\end{EESegment}

\begin{EESegment}{EE-TGT-P043-002}
Let \NSi{NS-F-P043-004}{\nu_k} be the norm of multiplication by
\NSi{NS-F-P043-005}{B^{-1}} on
\NSi{NS-F-P043-006}{C^k([-1,1])}. As in Lemma~\ref{lem:4.7}, let
\NSi{NS-F-P043-007}{q_k} be the bilinear norm of
\NSi{NS-F-P043-008}{\mathcal Q:C^k\times C^k\to C^k}. Since
\NSi{NS-F-P043-009}{\lVert d_N\rVert_{C^k}\le D_k/N}, it is enough to choose
\end{EESegment}
\NSFormula{NS-F-P043-010}{display}{none}
\[
N\ge 8\max_{k=0,2}\nu_k^2q_kD_k
\]
\begin{EESegment}{EE-TGT-P043-003}
to get a smooth solution of \eqref{eq:4.42} with
\NSi{NS-F-P043-011}{\lVert c\rVert_{C^2}\le 2\nu_2D_2/N}. Because the bump shapes are fixed,
this solution also satisfies
\end{EESegment}
\NSFormula{NS-F-P043-012}{display}{none}
\[
\max_{r\le 1}\lVert \partial_X^r(u_c,e_c)\rVert_2\le C/N.
\]
\begin{EESegment}{EE-TGT-P043-004}
When the five correction integrands in \eqref{eq:4.42} are integrated only up to
\NSi{NS-F-P043-013}{X\in[Y_0,Y_1]}, the result is
\end{EESegment}
\NSFormula{NS-F-P043-014}{display}{none}
\[
\max_{0\le j\le 2}\ \sup_{[Y_0,Y_1]\times[-1,1]}
\left|\partial_\eta^j(\mathfrak m_f-\mathfrak m_N)\right|\le C/N.
\]
\begin{EESegment}{EE-TGT-P043-005}
The shear formulas and Lemma~\ref{lem:4.4}\textup{(ii)} show that increasing
\NSi{NS-F-P043-015}{N}, if needed, gives
\end{EESegment}
\NSFormula{NS-F-P043-016}{display}{none}
\[
E_f\ge \frac12\min_{[Y_0,Y_1]\times[-1,1]}E_0,
\qquad
\lVert a_f-a_N,b_f-b_N,p_f-p_N\rVert_0
\le C_{\mathrm{rep}}/N.
\]
\begin{EESegment}{EE-TGT-P043-006}
Here \NSi{NS-F-P043-017}{C_{\mathrm{rep}}} does not depend on
\NSi{NS-F-P043-018}{N}. Increase \NSi{NS-F-P043-019}{N} until the final triples stay in the
fixed compact neighborhood used to define \NSi{NS-F-P043-020}{C_\Psi}. The inequality in
\eqref{eq:4.41} then stays strict on the repair interval. More exactly, if
\NSi{NS-F-P043-021}{N\ge 4C_\Psi C_{\mathrm{rep}}/\mu_R}, then
\end{EESegment}
\NSFormula{NS-F-P043-022}{display}{none}
\[
\Psi_i(a_f,b_f,p_f)\ge \mu_R/2-C_\Psi C_{\mathrm{rep}}/N
\ge \mu_R/4>0.
\]
\begin{EESegment}{EE-TGT-P043-007}
This correction changes no profile and no cumulative integral before
\NSi{NS-F-P043-023}{Y_0}.
\end{EESegment}

\begin{EESegment}{EE-TGT-P043-008}
Equations~\eqref{eq:4.42} give
\NSi{NS-F-P043-024}{\mathfrak m_f(Y_1)=\mathfrak m_0(Y_1)}. All the changes to the profiles are
supported in \NSi{NS-F-P043-025}{(X_-,Y_1)}, so
\end{EESegment}
\NSFormula{NS-F-P043-026}{display}{4.43}
\begin{equation}
\label{eq:4.43}
\begin{aligned}
(U_f,E_f)&=(U_0,E_0) &&\text{for }X\le X_-\text{ and }X\ge Y_1,\\
\mathfrak m_f(X)&=\mathfrak m_0(X) &&\text{for }X\ge Y_1,\\
(\Pi_f,V_f,Q_{s,f},N_{s,f},p_f,\mathcal T_{0,f})
&=(\Pi_0^{\mathrm{prof}},V_0,Q_{s,0},N_{s,0},p_0,
\mathcal T_{0,0})
&&\text{for }X\ge Y_1.
\end{aligned}
\tag{4.43}
\end{equation}
\begin{EESegment}{EE-TGT-P043-009}
The last two equalities come from Lemma~\ref{lem:4.4}\textup{(i)}. Here
\NSi{NS-F-P043-027}{V_f\text{ and }V_0} are the radial-velocity profiles calculated from the final
and joined pairs. Every lower bound on \NSi{NS-F-P043-028}{N} used above is finite and depends
only on data that has already been fixed. One finite integer can therefore meet all of them at
once. From now on, write \NSi{NS-F-P043-029}{(U,E,\Pi)} for the final profiles and
\NSi{NS-F-P043-030}{\mathcal T_0} for their stress.
\end{EESegment}

\begin{EESegment}{EE-TGT-P043-010}
\textbf{Step 4: verify the six conclusions.} The final profiles satisfy the cone inequalities,
and beyond the correction interval their five cumulative integrals agree exactly with those of the
joined pair. These two facts now give the regularity, support, and endpoint estimates in the
theorem. For the regularity in part~\textup{(i)}, the changes in Steps 2--3 have compact support
above \NSi{NS-F-P043-031}{X_{\mathrm{an}}}. They therefore leave the analytic axis rectangle from
Proposition~\ref{prop:4.10} unchanged. Every other profile piece and every join between pieces is
smooth. Positivity of \NSi{NS-F-P043-032}{E} follows from positivity of
\NSi{NS-F-P043-033}{E_0}, the exponential in \eqref{eq:4.38}, and the bound for
\NSi{NS-F-P043-034}{E_f} on the repair interval. Near the axis,
\NSi{NS-F-P043-035}{E=\sqrt{2X}\phi/\mathsf C}. On that neighborhood,
\NSi{NS-F-P043-036}{F=\phi/\mathsf C, U, \Pi, V_0/X} have the regularity stated in
part~\textup{(i)}. This gives finite bounds for every fixed mixed derivative on every finite
rectangle.
\end{EESegment}

\begin{EESegment}{EE-TGT-P043-011}
The preserved pressure increment satisfies
\NSi{NS-F-P043-037}{C_p(\infty,\eta)=-\Pi_0(\eta)}. It follows that
\end{EESegment}
\NSFormula{NS-F-P043-038}{display}{none}
\[
\Pi(X,\eta)=\Pi_0(\eta)+C_p(X,\eta)
=C_p(X,\eta)-C_p(\infty,\eta)
=-\int_X^\infty\frac{E(x,\eta)^2}{2x}\,dx.
\]
\begin{EESegment}{EE-TGT-P043-012}
This proves both the pressure normalization in part~\textup{(i)} and the radial balance in
part~\textup{(ii)}. In particular,
\NSi{NS-F-P043-039}{\Pi_X=E^2/(2X)=F^2} extends smoothly to the axis.
Proposition~\ref{prop:4.2} gives the exact
\end{EESegment}
\NSPage{044}%
\begin{EESegment}{EE-TGT-P044-001}
tangential residual identities. The first edge has not changed, so the stress is zero for
\NSi{NS-F-P044-001}{X\le X_a}, and \eqref{eq:4.13} holds there. By \eqref{eq:4.43}, the outer
stress has not changed either, and it is zero for
\NSi{NS-F-P044-002}{X\ge X_b}.
\end{EESegment}

\begin{EESegment}{EE-TGT-P044-002}
The cone condition holds when \NSi{NS-F-P044-003}{X_a<X<X_b}. It was kept near both edges and
beyond \NSi{NS-F-P044-004}{Y_1}, and Steps 2--3 proved it on
\NSi{NS-F-P044-005}{[X_-,Y_1]}. In particular, \eqref{eq:4.23} gives
\end{EESegment}
\NSFormula{NS-F-P044-006}{display}{none}
\[
\mathcal T_{0,\theta}+t_s\mathcal T_{0,z}=F(P_c-v_s)>0.
\]
\begin{EESegment}{EE-TGT-P044-003}
The stress is therefore nonzero at every interior point, which finishes part~\textup{(ii)}.
\end{EESegment}

\begin{EESegment}{EE-TGT-P044-004}
Next come the uniform claims in part~\textup{(iii)}. Set
\NSi{NS-F-P044-007}{y_a=\log(X/X_a),\ y_b=\log(X_b/X)}. The factorizations
\eqref{eq:4.33} and \eqref{eq:4.32} extend
\NSi{NS-F-P044-008}{n=\mathcal T_0/|\mathcal T_0|} to the edges as follows:
\end{EESegment}
\NSFormula{NS-F-P044-009}{display}{none}
\[
n=\frac{B_a}{|B_a|}\quad\text{near }X_a,
\qquad
n=\frac{(b_\theta,y_b^6b_z)}{\sqrt{b_\theta^2+y_b^{12}b_z^2}}
\quad\text{near }X_b.
\]
\begin{EESegment}{EE-TGT-P044-005}
Their denominators stay positive on smaller collars. Since
\NSi{NS-F-P044-010}{B_a(0,\eta)=F(X_a,\eta)\mathfrak s(X_a,\eta)},
\end{EESegment}
\NSFormula{NS-F-P044-011}{display}{none}
\[
n(X_a,\eta)=\frac{(1,t_s)}{\sqrt{1+t_s^2}},
\qquad n(X_b,\eta)=(1,0),
\qquad t_s(X_b,\eta)=0.
\]
\begin{EESegment}{EE-TGT-P044-006}
The smooth functions \NSi{NS-F-P044-012}{F,a,v_s-2} are positive in the interior. The endpoint
bounds from Lemma~\ref{lem:4.9} and Proposition~\ref{prop:4.10} make them positive at
\NSi{NS-F-P044-013}{X_a,X_b} as well. Each one therefore has a positive minimum on the closed
annulus.
\end{EESegment}

\begin{EESegment}{EE-TGT-P044-007}
On that rectangle, define
\end{EESegment}
\NSFormula{NS-F-P044-014}{display}{none}
\[
A_n=n_\theta+t_sn_z,
\qquad B_n=n_z-t_sn_\theta.
\]
\begin{EESegment}{EE-TGT-P044-008}
Inside the annulus, \eqref{eq:4.23} gives
\end{EESegment}
\NSFormula{NS-F-P044-015}{display}{none}
\[
A_n=\frac{P_c-v_s}{|p_s-\mathfrak s|}>0.
\]
\begin{EESegment}{EE-TGT-P044-009}
At both edges, \NSi{NS-F-P044-016}{B_n=0}. Also,
\NSi{NS-F-P044-017}{A_n=\sqrt{1+t_s^2}} at \NSi{NS-F-P044-018}{X_a}, while
\NSi{NS-F-P044-019}{A_n=1} at \NSi{NS-F-P044-020}{X_b}. It follows that
\NSi{NS-F-P044-021}{A_n>0} everywhere on the closed rectangle, so define
\end{EESegment}
\NSFormula{NS-F-P044-022}{display}{none}
\[
G_n=2-(v_s-2)B_n^2/A_n^2.
\]
\begin{EESegment}{EE-TGT-P044-010}
Equation~\eqref{eq:4.22} gives
\end{EESegment}
\NSFormula{NS-F-P044-023}{display}{none}
\[
G_n=2-(v_s-2)J_c^2/(P_c-v_s)^2>0
\]
\begin{EESegment}{EE-TGT-P044-011}
in the interior, while \NSi{NS-F-P044-024}{G_n=2} at both edges. The functions
\NSi{NS-F-P044-025}{A_n,G_n} are therefore smooth and positive everywhere. Taking
\end{EESegment}
\NSFormula{NS-F-P044-026}{display}{none}
\[
\kappa=\min\left\{1,
\min_{[X_a,X_b]\times[-1,1]}A_n,
\min_{[X_a,X_b]\times[-1,1]}G_n\right\}>0
\]
\begin{EESegment}{EE-TGT-P044-012}
gives \eqref{eq:4.26}, with \NSi{NS-F-P044-027}{\kappa<2}. This proves part~\textup{(iii)},
including the directions at the endpoints.
\end{EESegment}

\begin{EESegment}{EE-TGT-P044-013}
For part~\textup{(iv)}, let \NSi{NS-F-P044-028}{c_a=t_1^2>0} be the inner exponential constant
from \eqref{eq:4.33}, and put
\end{EESegment}
\NSFormula{NS-F-P044-029}{display}{none}
\[
\zeta(X)=\exp(-c_a/y_a^2-4/y_b^2)\quad(X_a<X<X_b),
\qquad \zeta=0\quad\text{otherwise}.
\]
\begin{EESegment}{EE-TGT-P044-014}
This function is smooth, and it is flat at both edges. On a fixed inner collar, the ratio
\NSi{NS-F-P044-030}{\zeta/e^{-c_a/y_a^2}} stays between two positive constants. On a fixed outer
collar, the same is true of \NSi{NS-F-P044-031}{\zeta/e^{-4/y_b^2}}. The smooth, nonzero factors
in \eqref{eq:4.33} and \eqref{eq:4.32} then give
\end{EESegment}
\NSFormula{NS-F-P044-032}{display}{none}
\begin{align*}
|\mathcal T_0|&\ge c\zeta,& |\partial^\alpha\mathcal T_0|&\le C_\alpha\zeta y_a^{-m_\alpha}
&&\text{on the inner collar},\\
|\mathcal T_0|&\ge c\zeta,& |\partial^\alpha\mathcal T_0|&\le C_\alpha\zeta y_b^{-m_\alpha}
&&\text{on the outer collar}.
\end{align*}
\begin{EESegment}{EE-TGT-P044-015}
Each derivative of an exponential factor contributes only finitely many inverse powers of
\NSi{NS-F-P044-033}{y_a\text{ or }y_b}. On the compact interval between the collars,
\NSi{NS-F-P044-034}{\zeta>0}, \NSi{NS-F-P044-035}{|\mathcal T_0|>0}, and every fixed derivative
is bounded. With \NSi{NS-F-P044-036}{\delta=\min\{1,y_a,y_b\}}, these bounds combine to give
\eqref{eq:4.27}, which proves part~\textup{(iv)}.
\end{EESegment}
\NSPage{045}%
\begin{EESegment}{EE-TGT-P045-001}
Finally, both exact connections leave the total moment identities in part~\textup{(v)} unchanged.
For the last correction, this follows directly from
\end{EESegment}
\NSFormula{NS-F-P045-001}{display}{none}
\begin{align*}
(M,J,S)_f(\infty)&=(M,J,S)_0(\infty)=(0,0,0),\\
\int_0^\infty(H_f-H_0)\,dX&=I_f(Y_1)-I_0(Y_1)=0.
\end{align*}
\begin{EESegment}{EE-TGT-P045-002}
Together with Step 1, these equalities give \eqref{eq:4.28}. The exterior pair is still the heat
profile from Lemma~\ref{lem:4.8}, with the same
\NSi{NS-F-P045-002}{c_\infty}. The identity
\end{EESegment}
\NSFormula{NS-F-P045-003}{display}{none}
\[
q^{-A}c_\infty X^{-A}\mathcal H(2d/X)
=c_\infty s^{-A}\mathcal H(2\tau/s)
\]
\begin{EESegment}{EE-TGT-P045-003}
proves the formula for the physical exterior. Since
\NSi{NS-F-P045-004}{Y_1<\sup\mathcal I_4<X_v}, equation \eqref{eq:4.43} keeps
\NSi{NS-F-P045-005}{U=M=J=0} for
\NSi{NS-F-P045-006}{X\ge X_v}, with the same \NSi{NS-F-P045-007}{X_v} as in
Lemma~\ref{lem:4.8}. Equation~\eqref{eq:4.7} then gives
\end{EESegment}
\NSFormula{NS-F-P045-008}{display}{none}
\[
V_0=\frac{2\eta XU-2D\eta M-dM_\eta}{L}=0.
\]
\begin{EESegment}{EE-TGT-P045-004}
The radius \NSi{NS-F-P045-009}{X_v} comes before the last collar and after all four reserved
intervals. This proves part~\textup{(v)}. For part~\textup{(vi)}, the heat compensation used
\NSi{NS-F-P045-010}{\mathcal I_2}, while Step 3 used only
\NSi{NS-F-P045-011}{\mathcal I_1}. The intervals
\NSi{NS-F-P045-012}{\mathcal I_{\mathrm{pos}}=\mathcal I_3} and
\NSi{NS-F-P045-013}{\mathcal I_{\mathrm{mean}}=\mathcal I_4} therefore keep the exact profile
\eqref{eq:4.30}, proving part~\textup{(vi)}. Every choice was made in profile coordinates before
the physical \NSi{NS-F-P045-014}{q} or any later correction stage was introduced, as required.
\end{EESegment}
\end{proof}

\begin{EESegment}{EE-TGT-P045-005}
The profiles now produce the leading field from Section~\ref{sec:3}, with its stress supported in
the prescribed annulus. The next section corrects the higher-order residual terms found after
\eqref{eq:4.7}.
\end{EESegment}

\begin{EESegment}{EE-TGT-P045-006}
\section{Correcting the base flow to every order}
\label{sec:5}
\end{EESegment}

\begin{EESegment}{EE-TGT-P045-007}
Near the axis, the leading profiles
\NSi{NS-F-P045-015}{(E,U,V_0,\Pi)} from Theorem~\ref{thm:4.6} satisfy the balances in
Equations~\eqref{eq:4.7} and \eqref{eq:4.13}. Axial viscosity and the remaining radial terms still
appear in the residual, however. Add axisymmetric corrections to obtain the background
\NSi{NS-F-P045-016}{(u_B,p_B)} described in Section~\ref{sec:3} and
Proposition~\ref{prop:5.5}. The velocity \NSi{NS-F-P045-017}{u_B} is exactly divergence-free.
Its momentum residual is the negative cylindrical divergence of the physical annular stress
\NSi{NS-F-P045-018}{\mathcal T_{\mathrm{phys}}}, as in \eqref{eq:5.41}, plus an error
\NSi{NS-F-P045-019}{\mathcal E_B} that is bounded by every power of
\NSi{NS-F-P045-020}{q} as \NSi{NS-F-P045-021}{q\downarrow0}. The same bounds hold for every fixed
physical derivative on each compact range of profile coordinates. The annular stress becomes the
input to the wave construction in Section~\ref{sec:7}.
\end{EESegment}

\begin{EESegment}{EE-TGT-P045-008}
These remaining terms appear at successive powers of
\NSi{NS-F-P045-022}{q^{2h}} in Equations~\eqref{eq:5.3}--\eqref{eq:5.6}. At each order,
Lemma~\ref{lem:5.1} solves for a correction near the axis. Lemma~\ref{lem:5.2} then extends that
correction in the radial direction and imposes the five integral conditions that keep its stress
supported in the prescribed annulus. The induction in Proposition~\ref{prop:5.3} finishes the current
order before using its coefficients to build the next one.
\end{EESegment}

\begin{EESegment}{EE-TGT-P045-009}
The leading profiles and all constants used to construct them stay fixed throughout. The resulting
sequence of coefficients is the formal expansion \eqref{eq:5.1}; this does not assert that the
unmodified infinite series converges. After this sequence is built, Lemma~\ref{lem:5.4}
produces smooth fields for \NSi{NS-F-P045-023}{q>0} with the same asymptotic coefficients. The
proof of Proposition~\ref{prop:5.5} checks that this summation keeps incompressibility, the exterior
velocity, and the required bounds on the residual.
\end{EESegment}

\begin{EESegment}{EE-TGT-P045-010}
\subsection{Higher coefficients near the axis.}
\label{sec:5.1}
\end{EESegment}
\begin{EESegment}{EE-TGT-P045-011}
First solve the coefficient equations \eqref{eq:5.2}--\eqref{eq:5.6} near the axis, where the
stress must remain zero. For \NSi{NS-F-P045-024}{n\ge0}, set
\NSi{NS-F-P045-025}{\lambda_n=2nh}, \NSi{NS-F-P045-026}{E_0=E},
\NSi{NS-F-P045-027}{U_0=U}, and \NSi{NS-F-P045-028}{\Pi_0=\Pi}. Every coefficient profile is a
function of \NSi{NS-F-P045-029}{(X,\eta)}. Here the subscript zero now marks the order in the
expansion. In particular, \NSi{NS-F-P045-030}{\Pi_0(X,\eta)} is the whole leading pressure
profile, and its trace
\end{EESegment}
\NSPage{046}%
\begin{EESegment}{EE-TGT-P046-001}
\NSi{NS-F-P046-001}{\Pi_0(0,\eta)} is the axis pressure function defined in
\eqref{eq:4.31}. Use the similarity variables from \eqref{eq:4.1}, and write the physical
coefficients as
\end{EESegment}
\NSFormula{NS-F-P046-002}{display}{5.1}
\begin{equation}
\label{eq:5.1}
\begin{aligned}
u_{\theta,n}&=q^{-A+\lambda_n}E_n
=rq^{-A-1/2+\lambda_n}\phi_n/\mathsf C,
& E_n&=\sqrt{2X}\phi_n/\mathsf C,\\
u_{z,n}&=q^{-A+\lambda_n}U_n,
& ru_{r,n}&=q^{\lambda_n}V_n,
& p_n&=q^{-2A+\lambda_n}\Pi_n.
\end{aligned}
\tag{5.1}
\end{equation}
\begin{EESegment}{EE-TGT-P046-002}
Here \NSi{NS-F-P046-003}{\phi_n}, rather than \NSi{NS-F-P046-004}{E_n}, is the azimuthal
coefficient that is smooth at \NSi{NS-F-P046-005}{X=0}. The radial average is the operator
from \eqref{eq:4.6}, applied with \NSi{NS-F-P046-006}{\eta} held fixed:
\NSi{NS-F-P046-007}{\mathcal A_X(U_n)(X,\eta)
=X^{-1}\int_0^XU_n(x,\eta)\,dx}.
\end{EESegment}

\begin{EESegment}{EE-TGT-P046-003}
The gap \NSi{NS-F-P046-008}{2h} between successive powers in \eqref{eq:5.1} accounts for the two
lower-order effects identified after \eqref{eq:4.7}. Recall from \eqref{eq:4.1} that
\NSi{NS-F-P046-009}{A=\tfrac12+h\text{ and }D=\tfrac12-h}, so
\NSi{NS-F-P046-010}{1-2D=2A-1=2h}. For the first effect, the axial derivative identity in
\eqref{eq:4.2} contributes the factor
\NSi{NS-F-P046-011}{q^{-2D}=q^{-1}q^{2h}} when two axial derivatives are taken. The radial
relation \NSi{NS-F-P046-012}{X=r^2/(2q)} in \eqref{eq:4.1} contributes
\NSi{NS-F-P046-013}{q^{-1}} when two radial derivatives are taken. This is why axial viscosity
applied to order \NSi{NS-F-P046-014}{n-1} appears in the angular and axial equations at order
\NSi{NS-F-P046-015}{n}. These are the \NSi{NS-F-P046-016}{\text{order-}n-1} viscous terms in
Equations~\eqref{eq:5.3} and \eqref{eq:5.4} below.
\end{EESegment}

\begin{EESegment}{EE-TGT-P046-004}
For the second effect, recall that the leading radial balance
\NSi{NS-F-P046-017}{r\partial_rp^{(0)}=(u_\theta^{(0)})^2} is the balance behind the pressure
identity in \eqref{eq:4.7}. The physical powers in \eqref{eq:4.3} give both sides the power
\NSi{NS-F-P046-018}{q^{-2A}}. After the other terms in the radial momentum equation are
multiplied by \NSi{NS-F-P046-019}{r}, they begin at
\NSi{NS-F-P046-020}{q^{-1}=q^{-2A}q^{2h}}, as noted after \eqref{eq:4.7}. This accounts for the
shift by one order in the pressure equation \eqref{eq:5.5}. The remaining radial terms are
collected in \eqref{eq:5.6} and enter one order later.
\end{EESegment}

\begin{EESegment}{EE-TGT-P046-005}
To include derivatives of the powers of \NSi{NS-F-P046-021}{q}, recall the operators
\NSi{NS-F-P046-022}{T_a,Z_a} from \eqref{eq:4.2}, and abbreviate
\end{EESegment}
\NSFormula{NS-F-P046-023}{display}{none}
\[
T_{a,n}=T_{a+\lambda_n},
\qquad Z_{a,n}=Z_{a+\lambda_n},
\qquad Z_{a,n}^{[2]}=Z_{a+\lambda_n-D}Z_{a+\lambda_n}.
\]
\begin{EESegment}{EE-TGT-P046-006}
The order of composition in the last expression keeps track of the change in the power after the
first axial derivative. Any field with a negative index is taken to be zero. With the integration
constant fixed by regularity at the axis, incompressibility is exactly the pair of identities
\end{EESegment}
\NSFormula{NS-F-P046-024}{display}{5.2}
\begin{equation}
\label{eq:5.2}
\partial_XV_n=-Z_{-A,n}U_n,
\qquad
\frac{V_n}{X}
=\frac{2\eta U_n-2\eta(D+\lambda_n)\mathcal A_X(U_n)
-d\partial_\eta\mathcal A_X(U_n)}{L}.
\tag{5.2}
\end{equation}
\begin{EESegment}{EE-TGT-P046-007}
Fix \NSi{NS-F-P046-025}{n\ge1}. At this point every coefficient of order
\NSi{NS-F-P046-026}{j<n} is known, including its radial cutoffs and the integral corrections from
Lemma~\ref{lem:5.2}. On one common inner interval, solve for
\NSi{NS-F-P046-027}{\phi_n,U_n,\Pi_n}. Equation~\eqref{eq:5.2} then determines
\NSi{NS-F-P046-028}{V_n} from \NSi{NS-F-P046-029}{U_n}. The pressure coefficient
\NSi{NS-F-P046-030}{\Pi_n} has to be solved at the same time as
\NSi{NS-F-P046-031}{\phi_n,U_n}.
\end{EESegment}

\begin{EESegment}{EE-TGT-P046-008}
Set \NSi{NS-F-P046-032}{b=-A-\tfrac12} and \NSi{NS-F-P046-033}{c=-A}. In the system below,
primes mean \NSi{NS-F-P046-034}{\partial_X}, with
\NSi{NS-F-P046-035}{\eta} held fixed. The following equations cancel the momentum residual of
order \NSi{NS-F-P046-036}{n} on the inner interval:
\end{EESegment}
\NSFormula{NS-F-P046-037}{display}{5.3}
\begin{equation}
\label{eq:5.3}
2(X\phi_n''+2\phi_n')
=T_{b,n}\phi_n
+\sum_{i+j=n}\left\{V_i(\phi_j'+\phi_j/X)+U_iZ_{b,j}\phi_j\right\}
-Z_{b,n-1}^{[2]}\phi_{n-1},
\tag{5.3}
\end{equation}
\NSFormula{NS-F-P046-038}{display}{5.4}
\begin{equation}
\label{eq:5.4}
2(XU_n''+U_n')
=T_{c,n}U_n
+\sum_{i+j=n}\left\{V_iU_j'+U_iZ_{c,j}U_j\right\}
+Z_{-2A,n}\Pi_n-Z_{c,n-1}^{[2]}U_{n-1},
\tag{5.4}
\end{equation}
\NSFormula{NS-F-P046-039}{display}{5.5}
\begin{equation}
\label{eq:5.5}
\Pi_n'=\mathsf C^{-2}\sum_{i+j=n}\phi_i\phi_j-\frac{\Omega_{n-1}}{2X},
\tag{5.5}
\end{equation}
\NSFormula{NS-F-P046-040}{display}{5.6}
\begin{equation}
\label{eq:5.6}
\Omega_k=T_{0,k}V_k
+\sum_{i+j=k}\left\{V_i(V_j'-V_j/(2X))+U_iZ_{0,j}V_j\right\}
-2XV_k''-Z_{0,k-1}^{[2]}V_{k-1}.
\tag{5.6}
\end{equation}
\NSPage{047}%
\begin{EESegment}{EE-TGT-P047-001}
Every \NSi{NS-F-P047-001}{\Omega_{n-1},\phi_{n-1},\text{ or }U_{n-1}} that appears on the
right-hand sides uses coefficients already constructed. For example, the pressure equation can be
written as
\end{EESegment}
\NSFormula{NS-F-P047-002}{display}{none}
\[
\Pi_n'=2\mathsf C^{-2}\phi_0\phi_n
+\mathsf C^{-2}\sum_{i=1}^{n-1}\phi_i\phi_{n-i}
-\frac{\Omega_{n-1}}{2X}.
\]
\begin{EESegment}{EE-TGT-P047-002}
When \NSi{NS-F-P047-003}{n=1}, the sum is empty. Every term in that sum and
\NSi{NS-F-P047-004}{\Omega_{n-1}} is already known, while the first term is linear in the current
unknown. The same point applies to the other equations after their transport sums are split into
the cases \NSi{NS-F-P047-005}{(i,j)=(0,n),(n,0)} and
\NSi{NS-F-P047-006}{1\le i,j<n}. The system is therefore linear at every positive order. The
pressure equation will also be enforced when the current profiles are extended beyond the inner
interval.
\end{EESegment}

\begin{EESegment}{EE-TGT-P047-003}
Here is how the powers of \NSi{NS-F-P047-007}{q} and the terms from the cylindrical basis arise.
For \NSi{NS-F-P047-008}{rq^{b+\lambda_n}\phi_n/\mathsf C}, angular radial transport contains
\NSi{NS-F-P047-009}{u_r(\partial_r+r^{-1})u_\theta}, which gives
\NSi{NS-F-P047-010}{V_i(\partial_X\phi_j+\phi_j/X)}. The angular vector Laplacian gives
\NSi{NS-F-P047-011}{2(X\partial_{XX}\phi_n+2\partial_X\phi_n)} after the factor
\NSi{NS-F-P047-012}{rq^{b-1+\lambda_n}/\mathsf C} is removed. The corresponding axial scalar
Laplacian is \NSi{NS-F-P047-013}{2(X\partial_{XX}U_n+\partial_XU_n)}. Every transport product
has the same power because \NSi{NS-F-P047-014}{A+D=1}. Axial viscosity moves one step forward in
the \NSi{NS-F-P047-015}{q^{2h}} expansion because
\NSi{NS-F-P047-016}{1-2D=2h}. In the radial equation, multiplying by
\NSi{NS-F-P047-017}{r} puts pressure and centrifugal acceleration at the power
\NSi{NS-F-P047-018}{q^{-2A+\lambda_n}}. The other radial terms begin at
\NSi{NS-F-P047-019}{q^{-1+\lambda_k}}, so \NSi{NS-F-P047-020}{k=n-1}. Finally, applying the
radial operator \NSi{NS-F-P047-021}{\partial_{rr}+r^{-1}\partial_r-r^{-2}} to
\NSi{NS-F-P047-022}{V(X)/r} gives
\NSi{NS-F-P047-023}{2X\partial_{XX}V/(qr)}, which gives the last two terms in \eqref{eq:5.6}.
\end{EESegment}

\begin{EESegment}{EE-TGT-P047-004}
Equation~\eqref{eq:5.2} gives \NSi{NS-F-P047-024}{V_j=Xv_j} with a smooth function
\NSi{NS-F-P047-025}{v_j}. It follows that every term in
\NSi{NS-F-P047-026}{\Omega_k} is divisible by \NSi{NS-F-P047-027}{X}. For example,
\NSi{NS-F-P047-028}{V_i(\partial_XV_j-V_j/(2X))
=Xv_i(v_j/2+X\partial_Xv_j)}. Thus \NSi{NS-F-P047-029}{\Omega_k/X} is regular at the axis,
including after any fixed number of parameter derivatives.
\end{EESegment}

\begin{EESegment}{EE-TGT-P047-005}
The linear equations must be solvable on the same radial interval at every order. For that to work,
the recursion needs bounds on its \NSi{NS-F-P047-030}{\eta}-derivatives even as the coefficient
bounds grow. The analytic dependence of the leading inner profile gives those bounds.
\end{EESegment}

\begin{lemma}
\label{lem:5.1}
\begin{EESegment}{EE-TGT-P047-006}
There is an interval \NSi{NS-F-P047-031}{0\le\xi\le a}, where
\NSi{NS-F-P047-032}{\xi=\sqrt X}, that starts at the axis and extends into the inner collar from
Theorem~\ref{thm:4.6}\textup{(i)}. The stress cone inequalities hold there with one uniform
positive margin. On this interval, every positive order has exactly one solution of
Equations~\eqref{eq:5.2}--\eqref{eq:5.6} with
\NSi{NS-F-P047-033}{\phi_n(0,\eta)=U_n(0,\eta)=\Pi_n(0,\eta)=0}. The profiles are smooth in
\NSi{NS-F-P047-034}{X}. For each fixed radial derivative, they are holomorphic on a neighborhood
of \NSi{NS-F-P047-035}{[-1,1]} in the parameter \NSi{NS-F-P047-036}{\eta}, which means that they
have a complex derivative at every point of that neighborhood. The neighborhood and
its bounds may depend on the order and the radial derivative, but
\NSi{NS-F-P047-037}{a} does not. The uniqueness statement applies among solutions whose profiles
extend holomorphically in \NSi{NS-F-P047-038}{\eta} to a neighborhood of
\NSi{NS-F-P047-039}{[-1,1]}, with bounds uniform for
\NSi{NS-F-P047-040}{0\le\xi\le a}.
\end{EESegment}
\end{lemma}

\begin{EESegment}{EE-TGT-P047-007}
Because every positive-order coefficient is zero at the axis, the leading traces of
\NSi{NS-F-P047-041}{\phi,U,\text{ and }\Pi} stay unchanged.
\end{EESegment}

\begin{proof}
\begin{EESegment}{EE-TGT-P047-008}
\textbf{Step 1: write a linear system on the fixed inner interval.} The analytic axis solution and
its first collar, supplied by Proposition~\ref{prop:B.2} and Theorem~\ref{thm:4.6}, have the stated
parameter regularity on one fixed radial rectangle. Choose \NSi{NS-F-P047-042}{a} inside this
rectangle and beyond the shorter collar that will remain uncut. At order
\NSi{NS-F-P047-043}{n}, use the lower-order coefficients already constructed, including their
cutoffs and moment corrections. The corrections are supported to the right of this rectangle, and
the radial cutoffs keep their analytic dependence on the parameter. A small enough complex
neighborhood \NSi{NS-F-P047-044}{S_\rho=\{\eta\in\mathbb C:
\operatorname{dist}(\eta,[-1,1])<\rho\}} therefore contains no zero of
\NSi{NS-F-P047-045}{L} and gives bounded holomorphic coefficients and source terms everywhere on
\NSi{NS-F-P047-046}{0\le\xi\le a}.
\end{EESegment}

\begin{EESegment}{EE-TGT-P047-009}
Introduce the extra variable \NSi{NS-F-P047-047}{K_n} and the vector
\NSi{NS-F-P047-048}{W_n} so that the system can be written as
\end{EESegment}
\NSFormula{NS-F-P047-049}{display}{5.7}
\begin{equation}
\label{eq:5.7}
\begin{aligned}
K_n&=\mathcal A_X(U_n)-U_n,
&W_n&=(\phi_n,U_n,K_n,\Pi_n,\partial_\xi\phi_n,\partial_\xi U_n)^{\mathsf T},\\
\partial_\xi W_n+\xi^{-1}\operatorname{diag}(0,0,2,0,3,1)W_n
&=A_0W_n+A_1\partial_\eta W_n+f_n,
&W_n(0)&=0.
\end{aligned}
\tag{5.7}
\end{equation}
\NSPage{048}%
\begin{EESegment}{EE-TGT-P048-001}
The matrices \NSi{NS-F-P048-001}{A_0,A_1} are
\NSi{NS-F-P048-002}{6\times6} matrices of coefficients in
\NSi{NS-F-P048-003}{(\xi,\eta)}, with \NSi{NS-F-P048-004}{n} fixed. The known source
\NSi{NS-F-P048-005}{f_n} has six components and is built from the lower-order profiles. The
extra unknown \NSi{NS-F-P048-006}{K_n} replaces the radial average with a first-order equation,
so no radial integral of a current unknown is left unevaluated in the system. In fact,
\NSi{NS-F-P048-007}{\partial_\xi K_n+2K_n/\xi=-\partial_\xi U_n}, and the two second-order radial
operators become \NSi{NS-F-P048-008}{(\partial_{\xi\xi}+3\xi^{-1}\partial_\xi)/2} and
\NSi{NS-F-P048-009}{(\partial_{\xi\xi}+\xi^{-1}\partial_\xi)/2}. At every positive order, every
term that contains an unknown of order \NSi{NS-F-P048-010}{n} is linear. The pressure row is
\NSi{NS-F-P048-011}{\partial_\xi\Pi_n=4\xi\phi_0\phi_n/\mathsf C^2+
\text{known terms}}. Use the same row for
\NSi{NS-F-P048-012}{X\partial_X\Pi_n} in the axial equation. With this form, every apparent
division by \NSi{NS-F-P048-013}{X} is regular.
\end{EESegment}

\begin{EESegment}{EE-TGT-P048-002}
\textbf{Step 2: solve without shrinking the radial interval.} Iterating the system takes
derivatives in \NSi{NS-F-P048-014}{\eta}. First control how many of these derivatives can appear,
and only then estimate the repeated radial integrals. The coefficient of
\NSi{NS-F-P048-015}{\partial_\eta} has the following sparse block form. Set
\NSi{NS-F-P048-016}{H_c=D\eta+dU_0}. Its only entries that may be nonzero are
\end{EESegment}
\NSFormula{NS-F-P048-017}{display}{none}
\begin{align*}
(A_1)_{51}&=2H_c/L,
& (A_1)_{52}=(A_1)_{53}&=-2d(\phi_0+X\partial_X\phi_0)/L,\\
(A_1)_{62}&=2(H_c-dX\partial_XU_0)/L,
& (A_1)_{63}&=-2dX\partial_XU_0/L,
& (A_1)_{64}&=2d/L.
\end{align*}
\begin{EESegment}{EE-TGT-P048-003}
Thus \NSi{NS-F-P048-018}{A_1}, at every radius and after any number of parameter derivatives,
maps the first four coordinates into the last two and sends the last two coordinates to zero.
\end{EESegment}

\begin{EESegment}{EE-TGT-P048-004}
For \NSi{NS-F-P048-019}{c_i=(0,0,2,0,3,1)_i}, invert the diagonal part of \eqref{eq:5.7},
with zero data at the axis, by using the integral operator
\end{EESegment}
\NSFormula{NS-F-P048-020}{display}{none}
\[
(Gg)_i(\xi)=\int_0^\xi(s/\xi)^{c_i}g_i(s)\,ds,
\qquad \mathcal K=G(A_0+A_1\partial_\eta).
\]
\begin{EESegment}{EE-TGT-P048-005}
The kernels of this operator do not depend on \NSi{NS-F-P048-021}{\eta}, and they keep the two
coordinate subspaces just described separate. If \NSi{NS-F-P048-022}{D_0} is any diagonal kernel
that occurs between two factors, then
\end{EESegment}
\NSFormula{NS-F-P048-023}{display}{none}
\[
A_1(\xi)D_0A_1(s)=0,
\qquad A_1(\xi)D_0\partial_\eta A_1(s)=0.
\]
\begin{EESegment}{EE-TGT-P048-006}
So two neighboring derivative operators give zero even when the derivative on the left lands on a
variable matrix coefficient. Any nonzero product of \NSi{NS-F-P048-024}{k} factors can therefore
contain at most \NSi{NS-F-P048-025}{p_k=\lceil k/2\rceil} parameter derivatives.
\end{EESegment}

\begin{EESegment}{EE-TGT-P048-007}
Bound \NSi{NS-F-P048-026}{A_0,A_1,f_n} on the larger neighborhood
\NSi{NS-F-P048-027}{S_\rho}, uniformly across the whole radial interval. On
\NSi{NS-F-P048-028}{S_{\rho'}}, where \NSi{NS-F-P048-029}{0<\rho'<\rho}, Cauchy's estimate bounds
each parameter derivative by using the distance from the smaller complex neighborhood to the edge
of the larger one. The available loss in radius is
\NSi{NS-F-P048-030}{\Delta=\rho-\rho'}, and it is divided among the at most
\NSi{NS-F-P048-031}{p_k} derivatives. The diagonal kernels have norm at most one. The nested radial
integrals are ordered so that every new integration variable is no larger than the preceding one.
Their region is an ordered simplex, whose volume is
\NSi{NS-F-P048-032}{a^{k+1}/(k+1)!}. Adding the at most
\NSi{NS-F-P048-033}{2^k} compositions gives
\end{EESegment}
\NSFormula{NS-F-P048-034}{display}{5.8}
\begin{equation}
\label{eq:5.8}
\lVert\mathcal K^kGf_n\rVert_{S_{\rho'}}
\le \frac{C_n^{k+1}a^{k+1}}{(k+1)!}
\left(\max\{1,p_k/\Delta\}\right)^{p_k}.
\tag{5.8}
\end{equation}
\begin{EESegment}{EE-TGT-P048-008}
The \NSi{NS-F-P048-035}{k}th root of the right-hand side of \eqref{eq:5.8} tends to zero. The
Picard series \NSi{NS-F-P048-036}{W_n=\sum_{k\ge0}\mathcal K^kGf_n} therefore converges on
\NSi{NS-F-P048-037}{S_{\rho'}} for every finite \NSi{NS-F-P048-038}{a}, no matter how large
\NSi{NS-F-P048-039}{C_n} is. Apply the same bound after iterating the homogeneous equation for the
difference of two solutions; this proves uniqueness. In particular, the radial interval does not
depend on the size of the coefficient bounds at later orders.
\end{EESegment}

\begin{EESegment}{EE-TGT-P048-009}
\textbf{Step 3: recover smooth profiles at the axis.} Radial smoothness follows without
differentiating a singular expression:
\end{EESegment}
\NSFormula{NS-F-P048-040}{display}{none}
\[
(Gg)_i=\xi\int_0^1t^{c_i}g_i(t\xi)\,dt,
\qquad
\partial_\xi(Gg)_i=g_i(\xi)-c_i\int_0^1t^{c_i}g_i(t\xi)\,dt.
\]
\begin{EESegment}{EE-TGT-P048-010}
The second identity is continuous at zero. Repeating this argument, and using a smaller
neighborhood of the parameter when needed, gives every fixed radial derivative. Extend the integral
\end{EESegment}
